% Diapositiva 23
% David
%=============================================
% . Mostrar los datos utilizados para la estimación 
%=============================================
\begin{frame}{Datos Utilizados para la Estimación}

    \begin{columns}[T]
        %========================================
        % COLUMNA IZQUIERDA: Tabla de Sensores (Inputs)
        %========================================
        \column{0.55\textwidth}
        \begin{block}{Entradas: Coordenadas y Observaciones}
            \centering
            \scriptsize 
            \begin{tabular}{@{}lcccc@{}}
                \toprule
                \textbf{Sensor} & \textbf{$x_i$ [km]} & \textbf{$y_i$ [km]} & \textbf{$z_i$ [km]} & \textbf{$Az_i$ obs.} \\ \midrule
                S1  & -4.0 & -3.0 &  0.2 & 0.005219 \\
                S2  & -2.5 &  2.8 & -0.1 & 0.009811 \\
                S3  &  0.0 &  4.0 &  0.1 & 0.012014 \\
                S4  &  3.5 &  3.0 & -0.2 & 0.024570 \\
                S5  &  4.5 & -1.5 &  0.0 & 0.080580 \\
                S6  &  2.0 & -4.0 &  0.3 & 0.071733 \\
                S7  & -1.0 & -4.5 & -0.1 & 0.027704 \\
                S8  & -4.5 &  1.0 &  0.2 & 0.003556 \\
                \textbf{S9}  & \textbf{0.5} & \textbf{0.5} & \textbf{0.0} & \textbf{0.852201} \\
                S10 &  2.8 &  0.8 & -0.3 & 0.380182 \\
                S11 & -1.8 & -1.2 &  0.1 & 0.119785 \\
                S12 &  1.0 &  3.2 &  0.2 & 0.037284 \\ \bottomrule
            \end{tabular}
        \end{block}

        %========================================
        % COLUMNA DERECHA: Parámetros Clave
        %========================================
        \column{0.42\textwidth}
        \vspace{0.5cm}

        \begin{block}{Parámetros del Problema Inverso}
            \centering
            \scriptsize
            \begin{tabular}{@{}ll@{}}
                \toprule
                \textbf{Elemento} & \textbf{Valor [km]} \\ \midrule
                Fuente real $(x_{\text{real}},\ y_{\text{real}},\ z_{\text{real}})$ & $(1.2, -0.8, -1.1)$ \\ \vspace{0.1cm}
                Amplitud inicial & $A_0 = 10$ \\ \vspace{0.1cm}
                Red instrumental & $12$ sensores \\ \vspace{0.1cm}
                Dato de ajuste & $Az_i$ observado \\ \bottomrule
            \end{tabular}
        \end{block}

        \vspace{0.5cm}
        
        \begin{block}{Objetivo}
            \scriptsize
            Utilizar el set de datos $Az_i$ para que el algoritmo recupere las coordenadas de la fuente real ocultas en la simulación.
        \end{block}

    \end{columns}

\end{frame}